\chapter{Introdução}
\label{introducao}

\epigraph{\textit{"O Esclarecimento é a saída do homem de sua menoridade, da qual ele próprio é culpado. A menoridade é a incapacidade de fazer uso de seu entendimento sem a direção de outro."}}{\textbf{--- Immanuel Kant}}

% ================================================================
\section{Contextualização}
% ================================================================

Os combustíveis fósseis correspondem, ainda hoje, a maior parte das emissões antrópicas de dióxido de carbono (\ce{CO2}) e de outros gases de efeito estufa; somente em 2019, a queima dessas fontes liberou algo em torno da ordem de $38 \pm 3{,}1$\,Gt\ce{CO2}  \cite{ipcc2023}. Embora a geração de energia elétrica e as indústrias concentrem a maior parcela absoluta dessas emissões, o setor de transportes apresenta uma particularidade que o distingue dos demais: entre 2010 e 2019, suas emissões cresceram a uma taxa aproximada de 2\% ao ano, num período em que outras áreas da economia já esboçavam, ainda que timidamente, trajetórias de redução \cite{ipcc2023}. Isso é, em parte, resultado das barreiras técnicas que ainda limitam a eletrificação direta de segmentos como o transporte rodoviário de cargas pesadas, a navegação e a aviação, o que torna necessária a busca por rotas complementares de descarbonização: biocombustíveis, hidrogênio de baixa emissão e combustíveis sintéticos figuram entre as alternativas com maior grau de maturidade tecnológica \cite{ipcc2023}.

Do ponto de vista da inserção imediata na matriz energética, o biodiesel se destaca entre os biocombustíveis líquidos por uma razão bastante prática: a sua compatibilidade com motores de ignição por compressão (também conhecidos como motores a Diesel) já existentes, sem que sejam necessárias modificações significativas nos equipamentos \cite{leung2010}. Quimicamente, trata-se de uma mistura de ésteres monoalquílicos de ácidos graxos de cadeia longa, produzidos pela reação de fontes lipídicas variadas (óleos vegetais, gorduras animais ou óleos residuais) com um álcool de cadeia curta, na presença de um catalisador \cite{dabdoub2009}. Essa compatibilidade, por sua vez, permite o aproveitamento da infraestrutura de distribuição e das frotas já em operação, o que reduz consideravelmente os custos associados à transição energética \cite{dabdoub2009}. Do ponto de vista das emissões locais, a substituição tende a diminuir a concentração de material particulado, monóxido de carbono e hidrocarbonetos não queimados nos gases de escape, além de eliminar os compostos sulfurados presentes no diesel fóssil \cite{encinar2016,verma2016}. Quanto ao balanço de carbono, o \ce{CO2} liberado na combustão corresponde ao carbono fixado previamente pela biomassa via fotossíntese, de modo que o ciclo tende à neutralidade quando as condições de manejo são adequadas \cite{ipcc2023}.

Essas vantagens, no entanto, não se realizam de forma independente da eficiência do processo produtivo. Análises de custo mostram que a matéria-prima responde por entre 70 e 88\% do preço final do biodiesel \cite{dabdoub2009}, de forma que qualquer ineficiência na etapa de conversão repercute diretamente na competitividade do produto frente ao diesel mineral. No Brasil, o estímulo à produção tem também dimensão institucional: o Programa Nacional de Produção e Uso do Biodiesel (PNPB) estipula percentuais obrigatórios de mistura ao diesel mineral, o que assegura ao mercado interno uma demanda estrutural pelo produto \cite{brasilpnpb2004}.

% ================================================================
\section{Processo de transesterificação e suas limitações}
% ================================================================

A transesterificação é o processo químico pelo qual um éster é convertido em outro mediante a troca do grupo alcóxido. Na produção de biodiesel, a reação ocorre entre triglicerídeos (principal classe de compostos dos óleos vegetais e gorduras animais) e um álcool de cadeia curta, tipicamente metanol ou etanol, na presença de um catalisador. O produto consiste em uma mistura de ésteres monoalquílicos de ácidos graxos (o biodiesel propriamente dito) e glicerol como coproduto\cite{dabdoub2009,verma2016}. A estequiometria global pode ser representada de forma simplificada como:

\begin{equation}
    \ce{TG + 3 ROH <=> Glicerol + 3 FAME}
    \label{eq:transesterificacao}
\end{equation}

\noindent em que TG representa o triglicerídeo e ROH o álcool empregado.
A reação global da Equação \ref{eq:transesterificacao} é, na realidade, composta por três etapas reversíveis consecutivas: o triglicerídeo é convertido em diglicerídeo, este em monoglicerídeo e, finalmente, em glicerol, com liberação de um éster alquílico a cada etapa \cite{roosta2016}.

A catálise homogênea alcalina (com \ce{NaOH}, \ce{KOH} ou metóxido de
sódio) é a rota predominante na produção industrial de biodiesel, em razão da elevada atividade catalítica, da disponibilidade comercial dos reagentes e das condições operacionais relativamente brandas, com temperaturas próximas a 60--65\,$^{\circ}$C e pressão atmosférica\cite{dabdoub2009,leung2010}. Essa rota, contudo, impõe restrições sobre a qualidade da matéria-prima: teores elevados de ácidos graxos livres e de água na carga favorecem a formação de sabão por saponificação, o que contamina o produto, dificulta a separação das fases e reduz o rendimento em ésteres\cite{verma2016}. A remoção desses contaminantes exige etapas adicionais de lavagem e purificação, com consumo significativo de água, energia e reagentes, o que eleva os custos operacionais de um processo cuja competitividade econômica já é fortemente dependente do preço da matéria-prima\cite{dabdoub2009}.

% ================================================================
\subsection{O problema da miscibilidade entre óleo e álcool}
% ================================================================

Um aspecto central da transesterificação alcalina reside na imiscibilidade entre os dois reagentes principais. Triglicerídeos são moléculas de elevada massa molar e caráter apolar, enquanto metanol e etanol são solventes polares com solubilidade limitada em óleos vegetais. Ao serem misturados, esses reagentes formam duas fases líquidas distintas, com área interfacial restrita entre elas\cite{roosta2016}. Como a reação de transesterificação ocorre na interface entre as fases, a transferência de massa passa a ser a etapa limitante do processo, e a velocidade global da reação é governada pela taxa com que os reagentes se encontram nessa região, e não pela atividade intrínseca do catalisador\cite{roosta2016}.

As consequências práticas dessa limitação são diretas. A primeira é a necessidade de agitação mecânica intensa para aumentar a área de contato entre as fases, o que eleva o consumo energético. A segunda é o emprego de álcool em excesso em relação à estequiometria: enquanto a reação requer 3 moles de álcool por mol de triglicerídeo (Equação\ref{eq:transesterificacao}), as condições industriais tipicamente utilizam razões molares álcool:óleo de 6:1 a mais de 10:1\cite{verma2016,ma2016}. Esse excesso visa deslocar o equilíbrio da reação reversível na direção dos produtos, mas sua origem está na baixa miscibilidade entre as fases, que impede o aproveitamento eficiente do álcool disponível. O álcool não reagido precisa ser recuperado por destilação ao final do processo, adicionando operações unitárias e custos ao ciclo produtivo\cite{dabdoub2009}.

Alternativas como o uso de condições supercríticas (temperaturas acima de 240\,$^{\circ}$C e pressões superiores a 8MPa) eliminam a barreira interfacial ao converter o sistema em uma fase única, mas a um custo energético e de infraestrutura que limita a aplicação industrial\cite{dabdoub2009}. Catalisadores heterogêneos, por sua vez, facilitam a separação do catalisador e evitam a formação de sabão, porém não resolvem o problema da imiscibilidade entre óleo e álcool\cite{leung2010}. Uma abordagem conceitualmente diferente consiste na adição de cossolventes ao meio reacional. Em vez de forçar amiscibilidade por meio de condições extremas de temperatura e pressão, o cossolvente modifica o ambiente químico de modo a tornar as duas fases parcial ou totalmente miscíveis em condições brandas\cite{roosta2016}. \citeonline{roosta2016} demonstraram que, em quantidade suficiente, o cossolvente pode gerar um sistema monofásico homogêneo, no qual a velocidade de reação é consideravelmente superior à do sistema bifásico convencional. Em quantidades menores, o efeito é o aumento da concentração mútua das espécies entre as fases, acelerando a reação mesmo sem eliminar completamente a heterogeneidade do sistema.

Solventes orgânicos convencionais como tetrahidrofurano e éter dietílico já foram avaliados nessa função\cite{encinar2016,roosta2016}, mas apresentam limitações práticas relevantes: volatilidade elevada, inflamabilidade e dificuldades na recuperação e no descarte. É nesse contexto, em que a eficácia da abordagem por cossolventes está demonstrada mas os solventes disponíveis impõem restrições operacionais significativas, que os líquidos iônicos se inserem como alternativa de interesse crescente.

% ================================================================
\section{Líquidos iônicos: definição, propriedades e relevância
         para o sistema reacional}
% ================================================================

Os líquidos iônicos (LIs) são sais compostos inteiramente por íons, tipicamente um cátion orgânico e um ânion orgânico ou inorgânico, que se encontram no estado líquido em temperaturas abaixo de 100\,$^{\circ}$C\cite{trindade2019}. Essa característica resulta da assimetria geométrica dos íons constituintes, que impõe impedimento estérico ao empacotamento ordenado necessário para a formação de um retículo cristalino estável\cite{alvarez2010}.

A natureza exclusivamente iônica desses compostos confere propriedades que os distinguem dos solventes orgânicos convencionais. A pressão de vapor dos LIs é negligenciável em condições normais de temperatura e pressão, o que elimina a emissão de compostos orgânicos voláteis durante o manuseio e a operação\cite{alvarez2010}. Além disso, a faixa de estabilidade líquida pode atingir 300\,$^{\circ}$C, permitindo que o produto da reação seja separado por destilação sem arraste ou perda do solvente iônico\cite{alvarez2010}. Outra característica relevante é a possibilidade de modular propriedades como hidrofobicidade, viscosidade e poder de solvatação por meio da escolha dos íons constituintes, o que confere aos LIs o caráter de \textit{solventes projetáveis}\cite{trindade2019,alvarez2010}.

Os LIs podem ser classificados em duas categorias conforme o mecanismo de formação. Os líquidos iônicos apróticos (LIAs) são sintetizados por reações de quaternização ou metátese, processos que envolvem múltiplas etapas, solventes auxiliares e purificação cuidadosa, o que resulta em custos de produção elevados\cite{troter2016,muhammad2015}. Os líquidos iônicos próticos (LIPs), por outro lado, são obtidos por transferência de próton de um ácido de Brønsted para uma base de Brønsted, em uma reação de neutralização direta que dispensa solventes e gera o produto em uma única etapa\cite{alvarez2010,trindade2019}. Essa simplicidade de síntese torna os LIPs consideravelmente mais baratos que os LIAs, o que é um fator decisivo quando se considera a aplicação em escala industrial. A menor toxicidade e a maior biodegradabilidade reportadas para vários LIPs reforçam o interesse por essa subclasse no contexto de processos sustentáveis\cite{trindade2019}.

Uma propriedade dos LIs particularmente relevante para o problema da miscibilidade óleo-álcool é a sua natureza anfifílica. Cátions orgânicos baseados em amônio, por exemplo, possuem uma região polar (a cabeça carregada positivamente) e uma ou mais cadeias alquílicas apolares. O ânion, conforme sua estrutura, pode também apresentar cauda hidrofóbica. Essa coexistência de regiões de polaridades distintas em uma mesma espécie iônica confere a diversos LIs um comportamento análogo ao de surfactantes\cite{alvarez2010}.

\citeonline{alvarez2010} investigaram uma família de LIPs baseados em cátions de $N$-metil-2-hidroxietilamônio com diferentes ânions carboxilatos, empregando espectroscopia DOSY-RMN e espalhamento de raios-X a baixo ângulo. Os resultados indicaram que, nos LIs contendo cadeia alquílica no ânion, as regiões hidrofóbicas se organizam espontaneamente em fases lamelares ou micelares, e os autores identificaram a existência de uma concentração micelar crítica (CMC) nesses compostos, fenômeno característico de sistemas anfifílicos. Para o sistema triglicerídeo/álcool da transesterificação, essa propriedade tem uma implicação prática: na interface entre as duas fases, espécies anfifílicas tendem a se posicionar com a cauda apolar voltada para a fase oleosa e a cabeça polar voltada para a fase alcoólica, reduzindo a tensão interfacial e aumentando a área de contato efetiva entre os reagentes\cite{trindade2019}.

Além do comportamento anfifílico, os LIPs apresentam uma característica estrutural que reforça sua adequação como cossolventes na transesterificação: a capacidade de formar redes extensas de ligações de hidrogênio. O próton transferido do ácido para a base durante a síntese do LIP gera um cátion com hidrogênio lábil (\ce{N-H}), capaz de atuar como doador e receptor de ligações de hidrogênio com álcoois como metanol e etanol\cite{trindade2019}. Essa interação, de natureza eletrostática e direcional, confere ao LIP uma afinidade particular pela fase alcoólica do sistema reacional. Quando combinada com a afinidade da cadeia alquílica apolar do ânion pela fase oleosa, o resultado é um composto com compatibilidade simultânea com ambas as fases, o que potencializa sua atuação na interface \cite{trindade2019,alvarez2010}.

A viabilidade do uso de LIPs como cossolventes é reforçada por considerações práticas. A baixa volatilidade desses compostos, associada à diferença de propriedades físicas em relação ao biodiesel e ao glicerol, permite a recuperação e a reutilização do LIP ao final do processo reacional\cite{troter2016}. Somam-se a isso a baixa toxicidade e a biodegradabilidade documentadas para LIPs à base de amônio, que representam uma vantagem ambiental concreta em relação aos solventes orgânicos convencionais empregados como cossolventes\cite{alvarez2010,trindade2019}.

% ================================================================
\section{Evidência experimental e lacuna na literatura}
% ================================================================

A eficácia de cossolventes na transesterificação alcalina está documentada para solventes orgânicos convencionais. \citeonline{encinar2016} demonstraram que a adição de éteres (éter dietílico, éter \textit{terc}-butilmetílico e tetrahidrofurano) ao sistema metanol/óleo de canola permitiu atingir teores de éster entre 97 e 98\% a 30\,$^{\circ}$C, temperatura consideravelmente inferior à do processo convencional. A explicação proposta pelos autores é que esses solventes possuem o equilíbrio de polaridades necessário para reduzir a tensão interfacial entre as fases\cite{encinar2016}. Todos os cossolventes testados, contudo, precisam ser completamente removidos do produto ao final da reação, em razão de sua volatilidade e toxicidade, o que impõe etapas adicionais de purificação e atenua parte do ganho cinético obtido\cite{encinar2016}.

Os LIPs oferecem uma alternativa que contorna essa limitação, e a evidência experimental disponível é promissora. \citeonline{trindade2019}avaliaram o hexanoato de dietilenotriamônio ([DETA][Hx]) como cossolvente na transesterificação alcalina de óleo de soja com metanol e observaram que a adição de apenas 3\% (m/m) do LIP promoveu um aumento de 8,14\% no rendimento mássico de biodiesel, atingindo cerca de 96\% em 60 minutos de reação. Os autores atribuem o resultado à baixa tensão superficial do [DETA][Hx] (24,3 mN/m a 60\,$^{\circ}$C), que intensifica a transferência de massa entre as fases \cite{trindade2019}.\citeonline{gomez2019}, do mesmo grupo de pesquisa, obtiveram resultado análogo com o hexanoato de 2-hidroxietilamônio ([2-HEA][Hx]): sob condições semelhantes (razão metanol/óleo 6:1, 1,5\% m/m de \ce{KOH}, 60\,$^{\circ}$C, 30min), a adição de 2\% (m/m) do LIP elevou o rendimento de 89,33\% para 92,72\%\cite{gomez2019}. Em ambos os casos, as concentrações necessárias são da ordem de 2 a 3\% (m/m), o que contrasta com o uso de LIs como catalisadores, função na qual concentrações muito mais elevadas são requeridas e comprometem a viabilidade econômica do processo\cite{muhammad2015,troter2016}.

A distinção entre esses dois modos de uso é relevante. Quando empregado como cossolvente, o LIP não substitui o catalisador alcalino convencional; sua função é remover a barreira de transferência de massa que limita a eficiência do catalisador já presente. Trata-se, portanto, de funções independentes e complementares, e a quantidade de LIP necessária para exercer o efeito cossolvente é substancialmente menor do que a requerida para uma ação catalítica\cite{trindade2019}.

Apesar desses resultados experimentais, a literatura ainda carece de ferramentas preditivas para orientar a seleção de LIs com potencial cossolvente. Os estudos disponíveis consistem em investigações experimentais de sistemas individuais, e a revisão bibliográfica conduzida neste trabalho não identificou modelos termodinâmicos validados para a predição do equilíbrio líquido-líquido (ELL) do sistema ternário \{triglicerídeo + álcool + líquido iônico prótico\} a partir de propriedades moleculares. Essa lacuna é particularmente crítica para modelos de contribuição de grupos, como o UNIFAC-LL, cuja aplicação depende da disponibilidade prévia de parâmetros de interação binários entre os grupos funcionais do sistema; para líquidos iônicos próticos, cujos grupos característicos (cátions amônio substituídos, ânions carboxilato) raramente constam das tabelas de parâmetros publicadas, essa dependência limita severamente o alcance preditivo do método. Modelos de natureza \textit{a priori}, como o COSMO-SAC, dispensam essa exigência ao calcular a atividade das espécies diretamente a partir de sua estrutura molecular, o que os torna candidatos naturais para sistemas nessa situação, mas sua adequação relativa frente aos métodos de contribuição de grupos, para o sistema ternário específico deste trabalho, ainda não havia sido estabelecida. Essa lacuna dupla (ausência de dados experimentais e ausência de comparação sistemática entre as duas classes de modelo) motiva o presente trabalho, que se propõe a avaliar comparativamente a capacidade preditiva do UNIFAC-LL e do COSMO-SAC para o ELL dos sistemas ternários \{triglicerídeo + álcool + líquido iônico prótico\}, tomando como líquidos iônicos próticos representativos o [2-HEA][Hx] e o [DETA][Hx].
% ================================================================
\section{Justificativa para a comparação entre UNIFAC-LL e COSMO-SAC}
% ================================================================

A predição do ELL de sistemas ternários \{triglicerídeo + álcool + cossolvente\} é o caminho natural para orientar, de forma racional, a seleção de LIPs com potencial cossolvente, mas a escolha do modelo termodinâmico empregado nessa predição não é trivial, e a literatura oferece duas abordagens de natureza distinta para esse problema.

A primeira abordagem é a dos modelos de coeficiente de atividade baseados em contribuição por grupos, como o UNIFAC-LL\cite{magnussen1981}, amplamente empregados na predição de equilíbrio de fases para sistemas orgânicos convencionais. Esses modelos calculam a atividade de cada espécie a partir de uma matriz de parâmetros de interação binários entre grupos funcionais, determinados por ajuste a dados experimentais de equilíbrio de fases. Essa dependência de dados prévios é, ao mesmo tempo, a força e a limitação do método: quando os parâmetros de um par de grupos já foram calibrados contra sistemas quimicamente semelhantes ao de interesse, a predição tende a ser confiável; quando não, o modelo precisa extrapolar para além do domínio em que foi validado, com confiabilidade decrescente. Para sistemas contendo líquidos iônicos, essa parametrização é particularmente limitada: as interações eletrostáticas de longo alcance e as redes de ligações de hidrogênio direcionais características dos LIs não foram contempladas na parametrização original desses modelos, e a cobertura para LIPs à base de amônio é escassa\cite{bahadur2015}.

A segunda abordagem é a dos modelos preditivos \textit{a priori}, entre os quais se destaca o COSMO-SAC (\textit{COnductor-like Screening MOdel -- Segment Activity Coefficient}), proposto por \citeonline{lin2002}. Diferentemente dos modelos de contribuição por grupos, o COSMO-SAC calcula os coeficientes de atividade a partir de perfis de densidade de carga superficial ($\sigma$-\textit{profiles}) obtidos por cálculos de química quântica sobre cada molécula isolada, sem requerer parâmetros ajustados a dados experimentais de misturas\cite{lin2002}. Essa natureza preditiva dispensa a existência prévia de um parâmetro de interação calibrado para o par de grupos específico do sistema, o que torna o modelo, em princípio, mais adequado a sistemas contendo espécies para as quais a contribuição por grupos carece de parametrização, como é o caso dos líquidos iônicos próticos estudados neste trabalho. A aplicabilidade da abordagem COSMO ao ELL de sistemas contendo LIs já foi demonstrada por \citeonline{dominguez2014}, que aplicaram o COSMO-RS a sistemas ternários com LIs de imidazólio, e por \citeonline{xu2022}, que empregaram o COSMO-SAC na predição de ELL de sistemas \{éster + álcool + LI\}.

Essa vantagem estrutural do COSMO-SAC, no entanto, não é absoluta nem dispensa verificação empírica: a qualidade de sua predição depende diretamente da geometria molecular e do perfil sigma que alimentam o cálculo, de modo que uma falha nessa etapa compromete o modelo de forma análoga a uma extrapolação malsucedida de parâmetros de grupo. Por essa razão, este trabalho não emprega o COSMO-SAC isoladamente, mas o aplica em paralelo ao UNIFAC-LL sobre o mesmo conjunto de sistemas ternários. Sendo assim, a comparação entre os dois modelos é aqui um instrumento, e não a finalidade: na ausência de dados experimentais que permitam verificar qualquer das duas predições, a convergência ou a divergência entre abordagens de fundamentos distintos é o único critério interno disponível para avaliar o quanto o comportamento previsto depende da escolha do modelo. Essa avaliação é retomada em detalhe no Capítulo~\ref{cap:resultados}.
% ================================================================
% ================================================================
\section{Objetivos}
\label{sec:objetivos}
 
O objetivo geral deste trabalho é avaliar, por meio de dois modelos termodinâmicos preditivos de naturezas distintas, o UNIFAC-LL e o COSMO-SAC, se a adição de líquido iônico prótico aumenta ou reduz a miscibilidade entre triglicerídeo e álcool no equilíbrio líquido-líquido a 298,15~K, combinando três triglicerídeos (trioleína, trilinoleína e trilinolenina), dois álcoois (etanol e metanol) e dois líquidos iônicos próticos (hexanoato de 2-hidroxietilamônio, [2-HEA][Hx], e hexanoato de dietilenotriamônio, [DETA][Hx]), verificando em que medida as duas abordagens convergem quanto ao comportamento de fases previsto.
 
Para atingir esse objetivo, os seguintes objetivos específicos são propostos:
 
\begin{itemize}
 
    \item Implementar as rotinas de cálculo de equilíbrio líquido-líquido (teste de estabilidade de fases, cálculo \textit{flash} e construção dos diagramas ternários) e verificar sua corretude de forma independente do modelo termodinâmico, bem como em conjunto com ele, contra sistemas de referência da literatura;
 
    \item Compor a matriz de parâmetros de interação do UNIFAC-LL para os sistemas estudados, suprindo por modelo preditivo de aprendizado de máquina as entradas ausentes na parametrização original para os grupos característicos dos líquidos iônicos próticos, e documentar a procedência e as limitações desse preenchimento;
 
    \item Definir a estratégia de representação molecular dos líquidos iônicos próticos no formalismo COSMO, incluindo o nível de teoria dos cálculos de química quântica e o tratamento do par iônico (cátion e ânion como entidades separadas \textit{versus} par iônico neutro), e gerar os $\sigma$-\textit{profiles} das espécies envolvidas, com verificação da qualidade geométrica de cada estrutura otimizada;
 
    \item Construir, para os dois modelos termodinâmicos, os diagramas ternários de equilíbrio líquido-líquido do conjunto completo de sistemas, e descrever o comportamento de fases previsto por cada um deles;
 
    \item Quantificar o efeito do líquido iônico prótico sobre a miscibilidade, a partir da evolução da largura da lacuna de imiscibilidade e da solubilidade do triglicerídeo na fase alcoólica em função da fração de cossolvente, tomando a aresta triglicerídeo-álcool como sistema de referência;
 
    \item Confrontar o comportamento obtido com os ganhos de rendimento reportados na literatura experimental, delimitando as diferenças de temperatura, de faixa de concentração e de grandeza medida que separam as duas situações;
 
    \item Interpretar os resultados à luz de mecanismos físico-químicos (impedimento estérico, competição pelo álcool e ligação de hidrogênio) e registrar as limitações comuns aos dois formalismos, na ausência de dados experimentais de equilíbrio de fases para os sistemas estudados.
 
\end{itemize}
% ================================================================
\section{Organização do trabalho}
\label{sec:organizacao}
% ================================================================

Este trabalho está organizado em cinco capítulos.

O Capítulo 1 apresenta a introdução, contextualizando a relevância do biodiesel como alternativa energética, os desafios do processo de transesterificação associados à imiscibilidade óleo-álcool e o potencial dos líquidos iônicos próticos como cossolventes. São apresentados também os objetivos do trabalho.

O Capítulo 2 desenvolve a fundamentação teórica, abordando os fundamentos termodinâmicos do equilíbrio líquido-líquido, a formulação do modelo COSMO-SAC e os princípios dos cálculos de química quântica em contínuo dielétrico (DFT/COSMO) sobre os quais o modelo se fundamenta.

O Capítulo 3 descreve a metodologia adotada, incluindo as fontes de dados experimentais utilizadas para avaliação do modelo, o nível de teoria empregado nos cálculos quânticos, a estratégia de representação dos líquidos iônicos próticos no formalismo COSMO e a operacionalização dos cálculos de ELL no \textit{software} JCOSMO.

O Capítulo 4 apresenta e discute os resultados. A ferramenta computacional é validada em duas etapas, isolando o motor de cálculo dos parâmetros e perfis específicos de cada modelo. Os diagramas ternários simulados pelo UNIFAC-LL e pelo COSMO-SAC são então apresentados por sistema e interpretados à luz de mecanismos físico-químicos, na ausência de dados experimentais para os sistemas ternários específicos deste trabalho.

O Capítulo 5 apresenta as conclusões, verificando o atendimento aos objetivos propostos e indicando direções para trabalhos futuros.

% ================================================================
\section{Declaração de uso de ferramentas de inteligência artificial}
\label{sec:declaracao_ia}
% ================================================================

Em atenção às boas práticas de integridade acadêmica, registra-se que ferramentas de inteligência artificial generativa foram empregadas como auxílio na elaboração deste trabalho. O uso restringiu-se à revisão textual e de estilo, à verificação da consistência interna da argumentação e das referências cruzadas, ao apoio na escrita e depuração de código e à identificação de lacunas e inconsistências no texto redigido pela autora.

A concepção do trabalho, ou seja, a definição dos sistemas estudados, os cálculos de equilíbrio líquido-líquido, a interpretação dos resultados e as decisões metodológicas e científicas neles apoiadas são de autoria própria, e a autora assume integral responsabilidade pelo conteúdo apresentado, incluindo a verificação de todas as informações, dados e referências aqui reunidos.